\documentclass[11pt,fontset=fandol]{ctexart}
\usepackage[a4paper,margin=23mm]{geometry}
\usepackage{amsmath,amssymb,mathtools,bm,mathrsfs,esint}
\usepackage{graphicx,longtable,booktabs,array,calc}
\usepackage[normalem]{ulem}
\usepackage{xcolor,hyperref}
\hypersetup{hidelinks}
\setcounter{secnumdepth}{-1}
\newcounter{none}
\providecommand{\tightlist}{\setlength{\itemsep}{0pt}\setlength{\parskip}{0pt}}
\providecommand{\passthrough}[1]{#1}
\usepackage{amssymb,mathrsfs,esint,bm,newunicodechar}
\allowdisplaybreaks
\setlength{\emergencystretch}{3em}
\providecommand{\llangle}{\langle\!\langle}
\providecommand{\rrangle}{\rangle\!\rangle}
\newunicodechar{∞}{\ensuremath{\infty}}
\newunicodechar{∈}{\ensuremath{\in}}
\newunicodechar{∑}{\ensuremath{\sum}}
\newunicodechar{∫}{\ensuremath{\int}}
\newunicodechar{≤}{\ensuremath{\leq}}
\newunicodechar{≥}{\ensuremath{\geq}}
\newunicodechar{≠}{\ensuremath{\neq}}
\newunicodechar{→}{\ensuremath{\to}}
\newunicodechar{⇒}{\ensuremath{\Rightarrow}}
\newunicodechar{α}{\ensuremath{\alpha}}
\newunicodechar{β}{\ensuremath{\beta}}
\newunicodechar{γ}{\ensuremath{\gamma}}
\newunicodechar{λ}{\ensuremath{\lambda}}
\newunicodechar{φ}{\ensuremath{\varphi}}
\newunicodechar{ε}{\ensuremath{\varepsilon}}
\newunicodechar{σ}{\ensuremath{\sigma}}
\newunicodechar{θ}{\ensuremath{\theta}}
\newunicodechar{δ}{\ensuremath{\delta}}
\newunicodechar{η}{\ensuremath{\eta}}
\newunicodechar{π}{\ensuremath{\pi}}
\newunicodechar{ξ}{\ensuremath{\xi}}
\newunicodechar{Σ}{\ensuremath{\Sigma}}
\newunicodechar{Δ}{\ensuremath{\Delta}}
\newunicodechar{∂}{\ensuremath{\partial}}
\newunicodechar{∀}{\ensuremath{\forall}}
\newunicodechar{∃}{\ensuremath{\exists}}
\newunicodechar{√}{\ensuremath{\surd}}
\newunicodechar{−}{\ensuremath{-}}
\newunicodechar{μ}{\ensuremath{\mu}}
\newunicodechar{ρ}{\ensuremath{\rho}}
\newunicodechar{τ}{\ensuremath{\tau}}
\newunicodechar{ω}{\ensuremath{\omega}}
\newunicodechar{Ω}{\ensuremath{\Omega}}
\newunicodechar{χ}{\ensuremath{\chi}}
\newunicodechar{ψ}{\ensuremath{\psi}}
\newunicodechar{Γ}{\ensuremath{\Gamma}}
\newunicodechar{Φ}{\ensuremath{\Phi}}
\newunicodechar{Λ}{\ensuremath{\Lambda}}
\newunicodechar{ℝ}{\ensuremath{\mathbb{R}}}
\newunicodechar{ℤ}{\ensuremath{\mathbb{Z}}}
\newunicodechar{ℕ}{\ensuremath{\mathbb{N}}}
\newunicodechar{△}{\ensuremath{\triangle}}
\newunicodechar{⊗}{\ensuremath{\otimes}}
\begin{document}
\section{2014年杭州电子科技大学601数学分析考研真题参考答案}\label{ux5e74ux676dux5ddeux7535ux5b50ux79d1ux6280ux5927ux5b66601ux6570ux5b66ux5206ux6790ux8003ux7814ux771fux9898ux53c2ux8003ux7b54ux6848}

来源：\texttt{HDU/601数学分析/2014/2014年杭州电子科技大学601数学分析考研真题参考答案.pdf}。原
PDF 页序：1, 2, 3, 4, 5, 6, 7, 8。

\section{杭州电子科技大学2014年数学分析真题解析}\label{ux676dux5ddeux7535ux5b50ux79d1ux6280ux5927ux5b662014ux5e74ux6570ux5b66ux5206ux6790ux771fux9898ux89e3ux6790}

\begin{quote}
助手整理：按同年真题与直接推导恢复水印污染、漏项和误识的数学公式。以下推导为重新整理的解答，非原参考答案逐字转录；涉及原题条件边界的地方明确说明。
\end{quote}

\subsection{一、填空题}\label{ux4e00ux586bux7a7aux9898}

\subsubsection{1. 极限}\label{ux6781ux9650}

令 \(t=\tan x\to0\)，则 \[
\boxed{\lim_{x\to0}(1+\tan x)^{\cot x}
=\lim_{t\to0}(1+t)^{1/t}=e.}
\]

\subsubsection{2.
分段函数连续}\label{ux5206ux6bb5ux51fdux6570ux8fdeux7eed}

原题函数为 \[
f(x)=\begin{cases}7e^x-\cos x,&x\leq0,\\\dfrac{\sin2x}{\tan ax},&x>0.\end{cases}
\] \(f(0)=6\)，右极限为 \(2/a\)，因此连续条件为 \(2/a=6\)，得
\(\boxed{a=1/3}\)。

\subsubsection{3.
变上限积分求导}\label{ux53d8ux4e0aux9650ux79efux5206ux6c42ux5bfc}

对 \(\int_0^xf(t)\,dt=x(1+\cos x)\) 求导，得
\(f(x)=1+\cos x-x\sin x\)，所以 \[
\boxed{f(\pi/2)=1-\pi/2.}
\]

\subsubsection{4. 定积分}\label{ux5b9aux79efux5206}

\(xe^{|x|}\) 是奇函数，\(|x|e^{|x|}\) 是偶函数，因此 \[
\boxed{\int_{-2}^2(|x|+x)e^{|x|}\,dx
=2\int_0^2xe^x\,dx=2[(x-1)e^x]_0^2=2(e^2+1).}
\]

\subsubsection{5.
复合函数二阶导数}\label{ux590dux5408ux51fdux6570ux4e8cux9636ux5bfcux6570}

\[
y'=e^{f(x)}f'(x),\qquad
\boxed{y''=e^{f(x)}\bigl([f'(x)]^2+f''(x)\bigr).}
\]

\subsubsection{6. 积分换序}\label{ux79efux5206ux6362ux5e8f}

由 \(x\leq y\leq2\sqrt x\)，等价地有 \(y^2/4\leq x\leq y\)，所以 \[
\boxed{\int_0^4dx\int_x^{2\sqrt x}f(x,y)\,dy
=\int_0^4dy\int_{y^2/4}^yf(x,y)\,dx.}
\]

\subsubsection{7.
幂级数的收敛域}\label{ux5e42ux7ea7ux6570ux7684ux6536ux655bux57df}

令 \(t=(x-3)/3\)，级数变为 \(\sum_{n\geq1}t^n/n\)。当 \(|t|<1\)
时绝对收敛，\(t=-1\) 时交错收敛，\(t=1\) 时发散，因此 \[
\boxed{\text{收敛域为 }[0,6).}
\]

\subsubsection{8.
隐函数二阶导数}\label{ux9690ux51fdux6570ux4e8cux9636ux5bfcux6570}

由 \(y-xe^y=1\)，写成 \(x=(y-1)e^{-y}\)。于是 \[
\frac{dx}{dy}=(2-y)e^{-y},\qquad
\frac{d^2x}{dy^2}=(y-3)e^{-y},
\] \[
\frac{d^2y}{dx^2}
=-\frac{d^2x/dy^2}{(dx/dy)^3}
=\frac{(3-y)e^{2y}}{(2-y)^3}.
\] 在 \(x=0\) 时 \(y=1\)，故 \(\boxed{y''(0)=2e^2}\)。

\subsection{二、计算题}\label{ux4e8cux8ba1ux7b97ux9898}

\subsubsection{1. 极限}\label{ux6781ux9650-1}

\[
\boxed{\lim_{x\to0}\frac1x\left(\frac1{\sin x}-\frac1{\tan x}\right)
=\lim_{x\to0}\frac{1-\cos x}{x\sin x}=\frac12.}
\]

\subsubsection{2. 不定积分}\label{ux4e0dux5b9aux79efux5206}

令 \(t=e^x\)，则 \[
\boxed{\int e^{e^x+x}\,dx=\int e^t\,dt=e^{e^x}+C.}
\]

\subsubsection{3.
幂级数求和及积分}\label{ux5e42ux7ea7ux6570ux6c42ux548cux53caux79efux5206}

在 \(|x|<1/3\) 内对几何级数求导，得到 \[
f(x)=\sum_{n=1}^{\infty}n3^{n-1}x^{n-1}
=\frac{d}{dx}\sum_{n=0}^{\infty}3^{n-1}x^n
=\boxed{\frac1{(1-3x)^2}}.
\] 因 \(1/8<1/3\)， \[
\boxed{\int_0^{1/8}f(x)\,dx
=\left.\frac1{3(1-3x)}\right|_0^{1/8}=\frac15.}
\]

\subsubsection{4.
右半圆曲线积分}\label{ux53f3ux534aux5706ux66f2ux7ebfux79efux5206}

将被积形式写成 \[
(e^x\sin y-my)\,dx+(e^x\cos y-m)\,dy
=d(e^x\sin y-my)-my\,dx.
\] 第一项积分为 \(\sin a-ma\)。当 \(a>0\) 时，路径从 \((0,0)\)
沿右半圆至 \((0,a)\)，与向下的直径组成正向边界，由 Green 公式 \[
\int_L-my\,dx=m\operatorname{area}(D)=\frac{\pi ma^2}{8}.
\] 所以在 \(a>0\) 时 \[
\boxed{I=\sin a-ma+\frac{\pi ma^2}{8}.}
\] 现有题面未显式写 \(a>0\)；若允许 \(a<0\)，方向相反，统一表达式应为
\(\sin a-ma+\pi ma|a|/8\)。可由参数
\(x=|a|\sin t/2\)、\(y=a(1-\cos t)/2\)、\(0\leq t\leq\pi\) 直接验证。

\subsection{三、综合题}\label{ux4e09ux7efcux5408ux9898}

\subsubsection{1.
原点处可微性与偏导连续性}\label{ux539fux70b9ux5904ux53efux5faeux6027ux4e0eux504fux5bfcux8fdeux7eedux6027}

\[
f(x,y)=\begin{cases}(x^2+y^2)\sin\dfrac1{x^2+y^2},&(x,y)\neq(0,0),\\0,&(x,y)=(0,0).\end{cases}
\] 令 \(r=\sqrt{x^2+y^2}\)，则 \[
\frac{|f(x,y)-f(0,0)|}{r}\leq r\longrightarrow0.
\] 所以 \(f\) 在原点可微，微分为零，且
\(f_x(0,0)=f_y(0,0)=0\)。在非原点处， \[
f_x=2x\sin\frac1{x^2+y^2}-\frac{2x}{x^2+y^2}\cos\frac1{x^2+y^2},
\] \[
f_y=2y\sin\frac1{x^2+y^2}-\frac{2y}{x^2+y^2}\cos\frac1{x^2+y^2}.
\] 取 \(x_n=(2\pi n)^{-1/2}\)、\(y_n=0\)，则
\(f_x(x_n,0)=-2/x_n\)，没有趋于零，故 \(f_x\) 在原点不连续。交换 \(x,y\)
可知 \(f_y\) 也在原点不连续。

\subsubsection{2.
函数项级数一致收敛}\label{ux51fdux6570ux9879ux7ea7ux6570ux4e00ux81f4ux6536ux655b}

原题的通项为 \(x/(1+n^4x^2)\)。由 \(2n^2|x|\leq1+n^4x^2\)， \[
\left|\frac{x}{1+n^4x^2}\right|\leq\frac1{2n^2}.
\] 控制级数收敛，因此原函数项级数在整个实轴上一致收敛。

\subsubsection{3.
球体积分的极限}\label{ux7403ux4f53ux79efux5206ux7684ux6781ux9650}

题设 \(f\) 在 \([0,1]\) 连续，\(f(0)=0\)、右导数 \(f'(0)=1\)，且 \[
F(t)=\iiint_{x^2+y^2+z^2\leq t^2}f(\sqrt{x^2+y^2+z^2})\,dx\,dy\,dz.
\] 球坐标变换给出 \(F(t)=4\pi\int_0^tr^2f(r)\,dr\)。因 \(f(r)=r+o(r)\)，
\[
F(t)=4\pi\int_0^t[r^3+o(r^3)]\,dr=\pi t^4+o(t^4).
\] 所以 \(\boxed{\lim_{t\to0^+}F(t)/t^4=\pi}\)。

\subsection{四、证明题}\label{ux56dbux8bc1ux660eux9898}

\subsubsection{1. 对数不等式}\label{ux5bf9ux6570ux4e0dux7b49ux5f0f}

对 \(u>0\)，\(\ln(1+u)=\int_0^u(1+t)^{-1}\,dt>u/(1+u)\)。取
\(u=1/x\)，即得 \(\ln(1+1/x)>1/(1+x)\)。

\subsubsection{2. Rolle 定理}\label{rolle-ux5b9aux7406}

令 \(F(x)=x[f(x)-f(a)]\)。由 \(f(a)=f(b)\)，\(F(a)=F(b)=0\)。故存在
\(\xi\in(a,b)\)，使 \[
0=F'(\xi)=f(\xi)-f(a)+\xi f'(\xi),
\] 即 \(\boxed{f(a)-f(\xi)=\xi f'(\xi)}\)。

\subsubsection{3.
连续但非一致连续}\label{ux8fdeux7eedux4f46ux975eux4e00ux81f4ux8fdeux7eed}

\(f(x)=\cos(1/x)\) 在 \((0,1)\) 上连续。取 \[
x_n=\frac1{2\pi n},\qquad y_n=\frac1{(2n+1)\pi}.
\] 两列均位于 \((0,1)\) 且 \(|x_n-y_n|\to0\)，但
\(|f(x_n)-f(y_n)|=2\)，故 \(f\) 在 \((0,1)\) 上不一致连续。

原答案附带的一般结论为：有限开区间上的连续函数一致连续，当且仅当两个端点的单侧极限都存在且有限。必要性由一致连续性和端点处的
Cauchy
准则得到；充分性通过在两个端点赋极限值，将函数连续延拓到闭区间，再用
Heine--Cantor 定理得到。相同证明适用于只缺一个端点的半开区间。

\subsubsection{4.
积分相等点的唯一性}\label{ux79efux5206ux76f8ux7b49ux70b9ux7684ux552fux4e00ux6027}

令 \[
F(t)=\int_a^tf(x)\,dx-\int_t^b\frac{dx}{f(x)}.
\] 因 \(f\) 连续且正，\(F\) 连续可导，且 \(F'(t)=f(t)+1/f(t)>0\)。又有
\[
F(a)=-\int_a^b\frac{dx}{f(x)}<0,\qquad F(b)=\int_a^bf(x)\,dx>0.
\] 由介值定理存在零点，由严格递增性知零点唯一，得到题目所需的唯一
\(\xi\in(a,b)\)。

\subsubsection{5.
截距和及原题的适用范围}\label{ux622aux8dddux548cux53caux539fux9898ux7684ux9002ux7528ux8303ux56f4}

对曲面 \(\sqrt x+\sqrt y+\sqrt z=\sqrt a\)，当所取点满足
\(x_0,y_0,z_0>0\) 时，其切平面为 \[
\frac{x-x_0}{2\sqrt{x_0}}+\frac{y-y_0}{2\sqrt{y_0}}+\frac{z-z_0}{2\sqrt{z_0}}=0,
\] 即 \[
\frac{x}{\sqrt{x_0}}+\frac{y}{\sqrt{y_0}}+\frac{z}{\sqrt{z_0}}=\sqrt a.
\] 三个截距依次为 \(\sqrt{ax_0},\sqrt{ay_0},\sqrt{az_0}\)，其和为 \[
\boxed{\sqrt a(\sqrt{x_0}+\sqrt{y_0}+\sqrt{z_0})=a.}
\] 原题``任意一点''的文字需要上述截距均有定义的条件。例如在边界点
\((0,a/4,a/4)\) 处，通常的三个非退化坐标轴截距并不存在，不能直接代入含
\(1/\sqrt{x_0}\)
的公式。这里保留原题问题并明确证明的适用范围，不把边界情形冒充为已满足条件。
\end{document}
